% ============================================================
% Capitolo 12 — Moduli e Skill
% ============================================================

\chapter{Moduli e Skill: caricare solo ciò che serve}
\label{ch:moduli-e-skill}

Il framework AI-DLC è composto da file che non vengono caricati tutti insieme all'inizio di ogni sessione. Vengono caricati selettivamente, in base alla fase e al tipo di task. Questo evita il ``context overload'' --- la situazione in cui l'agente ha così tante regole da non riuscire ad applicarne nessuna bene.

\section{La distinzione tra moduli e skill}
\label{sec:12-distinzione}

\textbf{I moduli} (\texttt{.ai-dlc/modules/NN\_*.md}) definiscono le regole operative del framework per ciascuna fase. Sono le istruzioni strutturali: come comportarsi, cosa verificare, quando fermarsi. Sono read-only.

\textbf{Le skill} (\texttt{.ai-dlc/modules/skills/SKILL\_*.md}) sono guide tematiche specializzate. Definiscono come affrontare un tipo specifico di lavoro.

\section{La mappa dei moduli}
\label{sec:12-mappa-moduli}

\begin{table}[htbp]
  \centering
  \begin{tabular}{lll}
    \toprule
    \textbf{File} & \textbf{Fase} & \textbf{Quando si carica} \\
    \midrule
    \texttt{00\_MODE.md} & Sempre & Regole per ogni Mode \\
    \texttt{01\_CORE\_RULES.md} & Sempre & Regole base, risk, halt, confidence tag \\
    \texttt{02\_DISCOVERY\_ANALYSIS.md} & 0--1 & Discovery e analisi requisiti \\
    \texttt{03\_DESIGN.md} & 2 & Architettura e design \\
    \texttt{04\_IMPLEMENTATION.md} & 3 & Implementazione e test \\
    \texttt{05\_VERIFICATION\_RELEASE.md} & 4--5 & QA e release \\
    \texttt{06\_OPS.md} & 6 & Operazioni e monitoring \\
    \texttt{07\_SPECIAL\_LANES.md} & Su richiesta & Hotfix, spike, parallel track \\
    \texttt{08\_PROMPT\_LIBRARY.md} & Su richiesta & Template di prompt \\
    \texttt{09\_CODEBASE\_ANALYSIS.md} & Su richiesta & Analisi codebase esistente \\
    \texttt{10\_DOCUMENTATION.md} & Su richiesta & Generazione documentazione \\
    \texttt{11\_BUGFIX\_PLAYBOOK.md} & Su richiesta & Debug strutturato \\
    \texttt{12\_OVERCONFIDENCE.md} & Sempre & Prevenzione overconfidence \\
    \texttt{13\_CONTENT\_VALIDATION.md} & Sempre & Validazione pre-scrittura \\
    \texttt{14\_STRUCTURED\_QUESTIONS.md} & Sempre & Domande strutturate \\
    \texttt{SEC\_CONSTRAINTS.md} & Su richiesta & Libreria vincoli SEC \\
    \texttt{PERF\_CONSTRAINTS.md} & Su richiesta & Libreria vincoli PERF \\
    \bottomrule
  \end{tabular}
  \caption{Mappa dei moduli AI-DLC}
  \label{tab:12-moduli}
\end{table}

\section{I moduli always-on: la triade della qualità}
\label{sec:12-always-on}

Oltre a \texttt{00\_MODE.md} e \texttt{01\_CORE\_RULES.md}, tre moduli sono sempre attivi. Non riguardano una fase specifica --- riguardano \emph{come} l'agente pensa e \emph{come} produce output.

\subsection*{\texttt{12\_OVERCONFIDENCE.md} --- Prevenzione dell'overconfidence}

Il problema più subdolo: un agente che assume requisiti non detti e procede con sicurezza su fondamenta instabili.

Il modulo impone:
\begin{itemize}
  \item \textbf{Red flag keywords:} se rispondi ``dipende'', ``forse'', ``standard'' senza specificare, l'agente non procede senza follow-up.
  \item \textbf{Anti-pattern espliciti:} non può completare una fase complessa senza domande, assumere requisiti non dichiarati, o scegliere lo stack senza conferma.
  \item \textbf{Soglie di confidenza:} $\geq 90\%$ procede dichiarando; 70--89\% chiede conferma; $< 70\%$ deve chiedere; $< 50\%$ si ferma.
\end{itemize}

\subsection*{\texttt{13\_CONTENT\_VALIDATION.md} --- Validazione pre-scrittura}

Ogni volta che l'agente scrive in un file, una checklist di validazione previene errori banali:
\begin{itemize}
  \item Code block bilanciati con language tag corretto
  \item Heading Markdown sequenziali, tabelle con colonne coerenti
  \item Diagrammi Mermaid/PlantUML con sintassi valida
  \item JSON senza virgole trailing, YAML con indentazione coerente
\end{itemize}

\subsection*{\texttt{14\_STRUCTURED\_QUESTIONS.md} --- Decisioni in file auditabili}

Le decisioni architetturali, di scope e di sicurezza vengono persistite in file \texttt{QS-<NN>-<desc>.md} invece di perdersi nella chat. Il workflow: l'agente identifica i punti di decisione, formula domande con opzioni A/B/C/D/E, salva il file, attende risposte. Ogni decisione è tracciabile (referenziata in ADR, \texttt{\_CONTEXT.md}, documenti di design).

\section{La mappa delle skill}
\label{sec:12-mappa-skill}

\begin{table}[htbp]
  \centering
  \begin{tabular}{ll}
    \toprule
    \textbf{File} & \textbf{Tipo di lavoro} \\
    \midrule
    \texttt{SKILL\_ANALYSIS.md} & Analisi requisiti, scenari d'uso \\
    \texttt{SKILL\_DESIGN.md} & Architettura, pattern, decisioni strutturali \\
    \texttt{SKILL\_API\_DESIGN.md} & API REST: naming, status code, paginazione \\
    \texttt{SKILL\_DATA\_ACCESS.md} & Schema DB, ORM, query, migration \\
    \texttt{SKILL\_SECURITY.md} & Revisione sicurezza, SEC-XX in pratica \\
    \texttt{SKILL\_TESTING.md} & Strategie di test, piramide, mocking \\
    \texttt{SKILL\_UI.md} & Componenti frontend, accessibilità \\
    \texttt{SKILL\_OPS.md} & CI/CD, monitoring, incident response \\
    \bottomrule
  \end{tabular}
  \caption{Skill disponibili nel framework}
  \label{tab:12-skill}
\end{table}

\section{Skill personalizzate di progetto}
\label{sec:12-skill-progetto}

Le skill del framework sono generiche. Il tuo progetto può avere esigenze specifiche in \texttt{.ai-dlc/project/skills/}:

\begin{lstlisting}[language=,caption={Skill di dominio per TaskFlow API}]
# SKILL — TaskFlow Domain

> Carica quando: lavori su task, progetti, assegnazioni.

## Regole di Dominio
- Un Task appartiene esattamente a un Project
- Transizioni stato: draft → open → in_progress → done | cancelled
- done e cancelled sono stati terminali
- completed_at viene impostato automaticamente su status → done

## Business Invariants
- Un utente non può essere assegnato a un task a cui non ha accesso
- Cancellazioni: soft delete (isDeleted = true, non physical delete)
\end{lstlisting}

\section{Il modulo \texttt{07\_SPECIAL\_LANES.md}}
\label{sec:12-special-lanes}

Copre tre casi fuori dal ciclo normale:

\textbf{Hotfix lane:} bug critico in produzione. Branch da main, fix minimale, review obbligatoria, rollback plan documentato prima del deploy.

\textbf{Spike lane:} esplorazione time-boxed (1--2 ore). Branch dedicato, obiettivo dichiarato, risultato in \texttt{PROGRESS.md} anche se lo spike fallisce.

\textbf{Parallel track:} feature su più layer (API + frontend). Tre sync point obbligatori: accordo sul contratto API, integrazione, E2E.

\begin{warnbox}
Non modificare mai \texttt{.ai-dlc/modules/}. Le personalizzazioni vanno in \texttt{.ai-dlc/project/} --- questo garantisce che il framework possa essere aggiornato senza perdere le tue regole.
\end{warnbox}

\section*{\LabelSummary}

\begin{itemize}
  \item I \textbf{moduli} sono le regole per ogni fase --- caricati automaticamente in base a \texttt{Phase}.
  \item Le \textbf{skill} sono guide tematiche --- caricate su richiesta in base al tipo di task.
  \item Skill personalizzate di progetto in \texttt{.ai-dlc/project/skills/} hanno priorità sulle skill del framework.
  \item I moduli del framework sono read-only: personalizza in \texttt{.ai-dlc/project/}.
\end{itemize}
